\input{prelude}
\usepackage{tikz}
\usetikzlibrary{arrows.meta,positioning}

\begin{document}

%%%%%%%%%%%%%%%%%%%%%%%%%%%
%%%%% ТИТУЛЬНЫЙ ЛИСТ %%%%%%
%%%%%%%%%%%%%%%%%%%%%%%%%%%
\begin{center} % включить выравнивание по центру
МИНИСТЕРСТВО ОБРАЗОВАНИЯ И НАУКИ РОССИЙСКОЙ \\
ФЕДЕРАЦИИ \\
Федеральное государственное автономное образовательное учреждение  \\
высшего образования\\ \textbf{<<Национальный исследовательский \\ Нижегородский государственный университет \\
им. Н.И. Лобачевского>> (ННГУ)}\\[1.5cm]%[4.5cm]
\textbf{Институт информационных технологий, математики и механики}\\[0.6cm]
\textbf{Кафедра: Прикладная математика}\\[0.6cm]
Направление подготовки: «Прикладная математика и информатика»\\
Профиль подготовки: «Прикладная математика и информатика (общий профиль)»\\[4.5cm]
%\textbf{Кафедра: Теории управления и динамики систем\\(ТУДС)}\\[0.8cm]
%Направление подготовки: \textbf{<<Математика и механика>>}\\
%направленность: \textbf{<<Дифференциальные уравнения, динамические системы и оптимальное управление>>}\\[2.2cm]

 \textbf{\large ВЫПУСКНАЯ КВАЛИФИКАЦИОННАЯ РАБОТА БАКАЛАВРА
} \\[0.3cm] % название работы, затем отступ 0,6см 
 
 на тему:\\
  \textbf{\large <<Параллельное программирование в задачах стохастической динамики>>}\\[5.5cm]
 % тема работы, затем отступ 3,7см
\begin{flushright}
 \begin{minipage}{0.52\textwidth} % начало маленькой врезки в половину ширины текста
 \begin{flushleft} % выровнять её содержимое по левому краю
  \textbf{Выполнил:} \\
 студент группы 3822Б1ПМ2 Морозов А.С.\\
  \underline{\hspace{4cm}} \textit{(подпись)} \\[0.3cm]
  \textbf{Проверил:}\\
    д.ф.-м.н., профессор Панкратов А.Л.\\
  \underline{\hspace{4cm}} \textit{(подпись)} \\[0.5cm]
 \end{flushleft} % конец выравнивания по левому краю
 \end{minipage} % конец врезки
\end{flushright}
 \vfill % заполнить всё доступное ниже пространство

  Нижний Новгород \\
 2026

 \thispagestyle{empty} % не нумеровать страницу

\end{center}

%%%%%%%%%%%%%%%%%%%%%%%%%%%%%%%%%%%%%
%%%%%%%% ОГЛАВЛЕНИЕ %%%%%%%%%%%%%%%%%
%%%%%%%%%%%%%%%%%%%%%%%%%%%%%%%%%%%%%
\renewcommand{\contentsname}{Содержание}
\newpage
\tableofcontents


%%%%%%%%%%%%%%%%%%%%%%%%%%%%%%%%%%%%%
%%%%%%%%% ОСНОВНАЯ ЧАСТЬ РАБОТЫ %%%%%
%%%%%%%%%%%%%%%%%%%%%%%%%%%%%%%%%%%%%
\newpage

\newpage
% --- НАЧАЛО ВВЕДЕНИЯ ---
\section*{Введение}
\addcontentsline{toc}{section}{Введение}

Сверхпроводимость представляет собой макроскопическое квантовое явление, при котором материал переходит в состояние с нулевым электрическим сопротивлением. Одним из наиболее значимых эффектов в физике сверхпроводников является эффект Джозефсона, описывающий туннелирование куперовских пар через тонкий слой диэлектрика, разделяющий два сверхпроводника. Джозефсоновские структуры лежат в основе современных высокотехнологичных устройств, включая сверхпроводниковые квантовые интерферометры (СКВИДы), высокочастотные генераторы и элементы квантовых вычислительных систем.

Динамика джозефсоновских переходов описывается нелинейными стохастическими дифференциальными уравнениями Ланжевена. Наличие шумовых флуктуаций приводит к возникновению сложных эффектов, при которых шум не только нарушает упорядоченное поведение системы, но и способен увеличивать время жизни метастабильных состояний. К таким эффектам относятся усиленная шумом стабильность (Noise Enhanced Stability, NES) и связанная с ней задержка переключения шумом (Noise Delayed Switching, NDS). Исследование подобных явлений требует проведения высокоточных численных экспериментов.

Для анализа стохастической динамики нелинейных осцилляторов второго порядка применяется численное моделирование стохастических дифференциальных уравнений. Получение статистически достоверных результатов требует моделирования большого ансамбля независимых стохастических траекторий. Основной вычислительной задачей является многократное численное интегрирование стохастических дифференциальных уравнений для большого числа независимых шумовых реализаций. Для получения достоверных результатов требуется интегрирование уравнений движения для десятков и сотен тысяч реализаций при малом временном шаге, что приводит к чрезвычайно высокой вычислительной сложности задачи.


\textbf{Актуальность темы} обусловлена необходимостью разработки высокопроизводительных параллельных алгоритмов для численного решения стохастических дифференциальных уравнений, моделирующих процессы в современных сверхпроводниковых устройствах.

\textbf{Объектом исследования} является стохастическая динамика нелинейного осциллятора второго порядка (джозефсоновского контакта) под воздействием шума и периодического сигнала.

\textbf{Предметом исследования} являются численные методы решения стохастических дифференциальных уравнений и методы их распараллеливания в многопоточных вычислительных средах.

\textbf{Цель работы} — разработка и исследование параллельных алгоритмов и программных средств численного моделирования стохастической динамики джозефсоновского перехода второго порядка, а также анализ индуцированных шумом эффектов.

В рамках работы реализован программный комплекс для моделирования ансамбля независимых стохастических траекторий и построения зависимостей среднего времени переключения (MST) от параметров шума и внешнего сигнала. Программа позволяет проводить параметрические свипы по $D$, $\alpha$, $A$ и $\omega$ и формировать статистические кривые на основе большого числа реализаций.

Для достижения поставленной цели были сформулированы следующие \textbf{задачи}:
\begin{enumerate}
    \item Описать физическую и математическую модель динамики фазы джозефсоновского контакта, эквивалентную движению частицы в периодическом потенциале.
    
    \item Разработать и реализовать численный алгоритм решения системы стохастических дифференциальных уравнений на основе стохастического метода Хюна (Heun).
    
    \item Разработать параллельную версию алгоритма для ансамблевого моделирования независимых стохастических траекторий.
    
    \item Исследовать особенности генерации псевдослучайных чисел в многопоточной вычислительной среде и обеспечить статистическую независимость шумовых реализаций.
    
    \item Провести вычислительные эксперименты и исследовать влияние интенсивности шума, коэффициента затухания и параметров внешнего сигнала на среднее время переключения системы.
    
    \item Исследовать эффекты усиленной шумом стабильности и задержки переключения шумом в системе второго порядка.
    
    \item Оценить вычислительную эффективность разработанного алгоритма и влияние распараллеливания на время проведения численных экспериментов.
\end{enumerate}

\newpage
% --- КОНЕЦ ВВЕДЕНИЯ ---

% --- НАЧАЛО ПЕРВОЙ ГЛАВЫ ---
\section{Физическая и математическая модель стохастической динамики}

\subsection{Эффект Джозефсона и динамика фазы}
Эффект Джозефсона, предсказанный Б. Джозефсоном в 1962 году, описывает туннелирование куперовских пар через барьер (тонкий слой диэлектрика), разделяющий два сверхпроводника. Волновые функции сверхпроводников по обе стороны барьера имеют определённые фазы, и разность этих фаз $\phi$ (или $x$) является ключевой макроскопической переменной системы. 

Динамика разности фаз джозефсоновского контакта может быть наглядно описана с помощью механической аналогии — как движение виртуальной броуновской частицы в наклонном периодическом потенциале («стиральной доске»). Потенциальная энергия такой частицы задается выражением:
\begin{equation}
    U(x) = 1 - \cos(x) - a x,
    \label{eq:potential}
\end{equation}
где первый член описывает периодический потенциал, связанный с взаимодействием куперовских пар, а второй член ($-ax$) учитывает влияние постоянного внешнего тока смещения, задающего наклон профиля потенциала.

Внешний ток существенно влияет на динамику фазы:
\begin{itemize}
    \item \textbf{Метастабильные состояния.} При токе $a < 1$ потенциальный профиль сохраняет локальные минимумы, разделённые барьерами. Частица (фаза) находится в одном из минимумов, совершая малые колебания. Это сверхпроводящее состояние контакта.
    \item \textbf{Исчезновение барьеров.} При достижении критического тока $a \ge 1$ локальные минимумы исчезают, частица начинает непрерывно скатываться по потенциалу. Возникает падение напряжения, и система переходит в резистивное состояние.
\end{itemize}

\subsection{Уравнение Ланжевена для осциллятора второго порядка}
Численные физические системы с диффузионными флуктуациями (джозефсоновские контакты, кольцевые лазеры, механические маятники) в рамках резистивной модели (RCSJ) описываются стохастическим дифференциальным уравнением Ланжевена. В отличие от сильно диссипативных моделей первого порядка, в данной работе исследуется полная динамика системы с учетом инертности (массы). 

Уравнение Ланжевена второго порядка для фазовой переменной $x(t)$ имеет вид:
\begin{equation}
    \ddot{x} + \alpha\dot{x} + \sin(x) = a + A \sin(\omega t) + g\cdot\xi(t),
    \label{eq:langevin_2nd_order}
\end{equation}
где:
\begin{itemize}
    \item $x(t)$ — координата системы (разность фаз);
    \item $\alpha$ — коэффициент затухания (вязкого трения), определяющий диссипацию энергии;
    \item $a$ — постоянная составляющая внешнего воздействия (ток смещения);
    \item $A, \omega$ — амплитуда и частота гармонического переменного сигнала;
    \item $\xi(t)$ — белый гауссов шум, моделирующий тепловые флуктуации;
    \item $g$ — интенсивность шумового воздействия (в коде часто $g = \sqrt{2D}$, где $D$ — температура или дисперсия).
\end{itemize}

Белый шум характеризуется нулевым математическим ожиданием $\langle\xi(t)\rangle = 0$ и дельта-корреляционной функцией $\langle\xi(t)\xi(t + \tau)\rangle = \delta(\tau)$. Наличие случайной силы позволяет системе преодолевать потенциальные барьеры даже при $a < 1$.

Для применения численных методов интегрирования уравнение (\ref{eq:langevin_2nd_order}) сводится к системе двух дифференциальных уравнений первого порядка:
\begin{equation}
\begin{cases}
    \frac{dx}{dt} = v \\
    \frac{dv}{dt} = a - \sin(x) - \alpha v + A \sin(\omega t) + g\cdot\xi(t)
\end{cases}
\label{eq:langevin_system}
\end{equation}

Наличие второй производной и переменной скорости $v$ существенно усложняет картину стохастических явлений по сравнению с моделями первого порядка. Инерционность добавляет эффекты низкочастотной фильтрации и кинетического (баллистического) пролета барьеров: накопив кинетическую энергию, частица может проскочить несколько минимумов подряд.

\subsection{Индуцированные шумом переходы и эффект задержки переключения}
Шумовые воздействия играют ключевую роль в динамике фазы вблизи метастабильных состояний. Если амплитуда детерминированного воздействия недостаточна для преодоления барьера, переход частицы из потенциальной ямы происходит исключительно за счет энергии флуктуаций (термоактивационные переходы).

Основной исследуемой характеристикой является среднее время переключения (Mean Switching Time, MST), соответствующее времени достижения частицей границы $x_{th} \approx \pi$. В отсутствие переменного сигнала ($A=0$) процесс подчиняется \textbf{закону Крамерса}: среднее время пребывания в яме спадает экспоненциально с ростом интенсивности шума $D$.

Однако при добавлении медленно меняющегося периодического сигнала возникает нетривиальный эффект конкуренции выталкивающей детерминированной силы и стохастического перемешивания. Этот эффект известен как \textbf{задержка переключения шумом} (Noise Delayed Switching, NDS) или усиленная шумом стабильность (Noise Enhanced Stability, NES). Эффект характеризуется следующим поведением системы в зависимости от уровня шума:
\begin{enumerate}
    \item \textbf{Малые шумы:} динамика близка к детерминированной. Частица покидает яму практически мгновенно при достижении подходящей фазы внешнего сигнала.
    \item \textbf{Промежуточные шумы:} наблюдается аномальный рост времени переключения. Шум «сбивает» траекторию частицы, забрасывая её обратно на дно потенциальной ямы в те моменты, когда эффективный барьер минимален. Возникает характерный «горб» на графике зависимости MST от $D$.
    \item \textbf{Сверхбольшие шумы:} термические флуктуации становятся доминирующими. Происходит разрушение синхронизации, и частица вылетает из ямы за счет термоактивации быстрее, чем меняется фаза внешнего сигнала (возврат к закону Крамерса).
\end{enumerate}

\subsection{Вычислительная сложность и необходимость распараллеливания}
Анализ стохастических эффектов, описанных выше, требует получения предельно гладких статистических кривых. Время одиночного перескока частицы $T$ является случайной величиной с большим разбросом. Для вычисления математического ожидания (MST) с приемлемой точностью необходимо провести усреднение по большому ансамблю реализаций. Обычно для получения одной точки на графике требуется смоделировать от $1 000$ до $10 000$ независимых траекторий.

Интегрирование системы СДУ (\ref{eq:langevin_system}) должно выполняться с малым временным шагом ($dt = 0.01$) для обеспечения устойчивости численной схемы. Таким образом, однократное построение графика зависимости MST от интенсивности шума (например, по 50 точкам шума) требует вычисления сотен миллионов итераций. При использовании последовательных алгоритмов на современном процессоре такой эксперимент может занимать от нескольких десятков минут до нескольких часов.

Поскольку вычисление каждой траектории не зависит от остальных (отсутствует информационная взаимозависимость по данным), данная задача относится к классу идеально параллельных (Embarrassingly Parallel). Это открывает возможности для применения многопоточности на уровне независимых траекторий и/или точек параметрического свипа (например, средствами стандартной библиотеки C++ и Qt), что позволяет эффективно задействовать многоядерные архитектуры и сократить время проведения численных экспериментов.
% --- КОНЕЦ ПЕРВОЙ ГЛАВЫ ---
\end{document}
